% !TEX root = ../main.tex

\unnumberedchapter{前言}

在正式开始本书之前，我们需要先回答一个问题：\textbf{什么是深度学习？}

在经典的机器学习方法中，计算机完成一个具体任务，通常需要人为设计特征。以图像分类为例，研究者需要根据经验提取边缘、纹理、形状等信息，再将这些特征输入分类器进行学习；对于文本处理，也需要预先构造词袋、词频等特征表示。这种方法在任务结构相对明确、数据规模有限的场景下具有较好的可解释性，但其性能高度依赖人工设计的特征质量。而随着数据规模不断扩大、数据形式日益复杂，人工设计特征逐渐成为机器学习中的瓶颈。

\textbf{深度学习（Deep Learning）}的核心思想，是利用具有多层结构的人工神经网络，让模型直接从数据中学习有效的特征表示。传统机器学习需要人工设计特征，再利用这些特征完成分类、预测等任务；而深度学习则可以通过端到端的训练，使模型在完成具体任务的同时自动学习适合当前任务的特征表示。训练过程中，模型利用反向传播和梯度下降不断调整网络参数，使模型逐步学习到更有效的特征表示。

深度学习中的“深度”二字，主要体现在网络层次的逐级堆叠。对于图像数据，网络可以从较低层次的边缘、纹理等局部模式开始，逐步组合出更复杂的形状、部件以及完整目标；对于文本数据，模型则可以从词或子词等基础表示出发，逐步学习上下文关系、句法结构以及更高层次的语义信息。因此，深度学习不仅是一类具体的模型结构，也是一种重要的\textbf{表征学习（Representation Learning）}范式。

随着计算能力、数据规模以及优化算法的发展，深度学习已广泛应用于计算机视觉、自然语言处理、语音识别和推荐系统等领域。近年来，以卷积神经网络、Transformer 等为代表的深度学习模型不断拓展应用范围，并推动了图像理解、序列建模以及语言生成等任务的发展。

本书默认读者具备以下基本数学基础：
\begin{itemize}
    \item \textbf{微积分：}掌握导数、偏导数、链式法则、梯度以及雅可比矩阵等基本概念；
    \item \textbf{线性代数：}掌握向量和矩阵的基本运算，以及行列式、秩、特征值与特征向量等基本概念；
    \item \textbf{概率论与数理统计：}掌握条件概率、常见概率分布、期望、方差、协方差以及极大似然估计等基本概念。
\end{itemize}

以上数学知识以能够理解和完成书中的相关推导为基本要求，不涉及过于深入或抽象的数学理论。对于书中涉及公式推导的部分，建议读者结合具体计算过程亲自动手演算，以加深对相关概念和计算过程的理解，而不是仅仅记忆结论。

\textbf{共勉。}

\vspace{1cm}
\begin{flushright}
\textbf{偷书的孔乙己} \\
\textit{\the\year 年于天津}
\end{flushright}

\cleardoublepage
